

\def\headerDeutsch{\selectlanguage{ngerman}
\textbf{Zulässige Hilfsmittel} sind ein blaues oder schwarzes, dauerhaftes, nicht radierbares Schreibwerkzeug.
\par
\textbf{Unzulässige Hilfsmittel} sind insbesondere Skripten, Folienkopien, Bücher, Mobiltelefone, handschriftliche Notizen, Headsets, Taschenrechner, Mikrophone, Smart Watches, vernetzte Geräte uä.
\par
\textbf{Hinweis}: Bei Betrugsversuch oder bei Nutzung unzulässiger Hilfsmittel wird die Arbeit 
nicht gewertet. 
\par
\textbf{Hinweis}: 
Antworten müssen direkt unterhalb der Frage in den Aufgabenbogen eingetragen werden.
Antworten, die an einer anderen Stelle im Aufgabenbogen eingetragen sind, können bei der Korrektur übersehen
werden und verzögern die Korrektur. Beantworten Sie die Frage so, daß der dafür 
verfügbare Platz ausreicht und dem Korrektor klar ist, was genau Ihre Antwort ist.
\par
\textbf{Hinweis}: In der Prüfung sollen Sie auch Ihre Fähigkeit unter Beweis stellen, 
klare und gut strukturierte Antworten zu geben. 
Gefragt sind nicht Stichworte oder Assoziationen, sondern Antworten auf die Fragen.
\par
\textbf{Hinweis}: Bitte schreiben Sie lesbar. Antworten, die wir nicht lesen können
können wir nicht bewerten. Sie können auf Deutsch oder auf Englisch antworten.
\par
\textbf{Hinweis}: Sie bekommen von uns auch ein leeres Blatt Papier für eigene
Notizen. Die 
Eintragungen in diesem Blatt dienen \textbf{nur für Ihre Vorbereitung} der Antwort und \textbf{werden nicht bewertet}. 
\par
\textbf{Hinweis}: Die Aufgabenblätter müssen vollständig abgegeben werden, \textbf{die
Notizen werden nicht abgegeben}.
\par
\textbf{Hinweis}: Bei einer Frage oder Toilettengang bitte ein Handzeichen geben. Die
Aufsicht kommt dann zu Ihnen an den Platz. 
\par
{\color{red}
\textbf{Erklärung}: Ich habe diese Hinweise verstanden und fühle mich gesundheitlich in
der Lage, die Prüfung anzutreten. Nach Antritt der Prüfung kann ich nicht mehr mit Hinweis
auf gesundheitliche Beeinträchtigungen die Prüfung abbrechen.
\par
}}


\def\headerEnglish{\selectlanguage{english}
\textbf{Permitted aids} are a blue or black, permanent, non-erasable writing tool.
\par
\textbf{Prohibited aids} include, in particular, scripts, photocopies of slides, books, mobile phones, handwritten notes, headsets, calculators, microphones, smartwatches, connected devices, etc.
\par
\textbf{Note}: In case of an attempt to cheat or the use of prohibited aids, the entire work will be invalidated.
\par
\textbf{Note}: Answers must be entered directly below the question into the answer sheet. Answers entered elsewhere in the sheet could be overlooked
and delay the process of grading.
Use the space given in the answer sheet and make it clear what exactly constitutes your answer.
\par
\textbf{Note}: In the exam, you must also demonstrate your ability to provide clear and well-structured answers. Keywords or associations are not requested, but responses to the posed questions are required.
\par
\textbf{Note}: Please write legibly. Answers that we cannot read cannot be evaluated.
You can answer in German or in English.
\par
\textbf{Note}: You will also receive a blank sheet of paper which is \textbf{only for your own preparation}.
The entries on this blank sheet \textbf{will not be evaluated}.
\par
\textbf{Note}: All answer sheets must be rendered at the end of the exam. The \textbf{note sheet is not collected}.
\par
\textbf{Note}: If you have a question or need to use the rest\-room, please raise your hand. 
The supervisor will then come to your seat.
\par
{\color{red}
\textbf{Declaration}: I have understood these instructions and feel physically capable of taking the exam. 
After starting the exam, I can no longer cancel the exam due to health reasons.
}
\par
}
